% !TEX root = ../main.tex

\chapter{进阶排版示例}
\label{chap:typesetting-examples}

本章展示算法、代码、矢量绘图和数学环境的写法。%
\footnote{本章算法、程序和数据用于演示排版，实际论文应替换为自己的研究内容。}
算法与定理的名称、编号采用 SJTUTeX 的适配；代码和图形由
\texttt{setup.tex} 中已加载的宏包提供。
全部示例可随论文直接编译；图表中的数值仅用于演示，不能作为实证结论。

\section{算法与伪代码}

\texttt{algorithm2e} 提供输入、输出、分支、循环和行号。
算法\ref{alg:binary-search} 展示二分查找，要求输入序列已经按升序排列。
每次迭代将候选区间缩小约一半，时间复杂度为 $O(\log n)$。
\texttt{\textbackslash label} 紧随标题书写，正文使用
\texttt{\textbackslash ref} 引用算法编号。

\begin{algorithm}[htbp]
  \caption{有序序列的二分查找}
  \label{alg:binary-search}
  \KwIn{升序序列 $A=(a_0,\ldots,a_{n-1})$，目标值 $x$}
  \KwOut{目标值的一个下标；未找到时返回 $-1$}
  $\ell \gets 0$，$r \gets n-1$\;
  \While{$\ell \le r$}{
    $m \gets \lfloor (\ell+r)/2 \rfloor$\;
    \eIf{$a_m=x$}{
      \KwRet $m$\;
    }{
      \eIf{$a_m<x$}{
        $\ell \gets m+1$\;
      }{
        $r \gets m-1$\;
      }
    }
  }
  \KwRet $-1$\;
\end{algorithm}

\section{程序代码}

\subsection{带行号和标题的 Python 代码}

代码\ref{lst:binary-search} 是算法\ref{alg:binary-search} 的 Python 实现。
\texttt{codeblock} 可设置语言、行号和标题，长行自动折行。

\begin{codeblock}[language=Python,float=htbp,numbers=left,numberstyle=\tiny\color{black!55},numbersep=8pt,caption={二分查找的 Python 实现},label={lst:binary-search},captionpos=b]
def binary_search(values, target):
    """Return a matching index, or -1 when absent."""
    left, right = 0, len(values) - 1
    while left <= right:
        middle = (left + right) // 2
        if values[middle] == target:
            return middle
        if values[middle] < target:
            left = middle + 1
        else:
            right = middle - 1
    return -1

assert binary_search([2, 5, 8, 13], 8) == 2
assert binary_search([2, 5, 8, 13], 7) == -1
\end{codeblock}

\subsection{SQL 与行内代码}

代码\ref{lst:sql-summary} 展示 SQL 聚合查询，假设
\lstinline!transactions! 表包含客户编号和交易金额两个字段。
行内代码可以使用 \texttt{\textbackslash lstinline}，例如
\lstinline!SUM(amount)!。
代码默认将美元符号等字符按原样显示；若需在注释中插入公式，可在单个代码块中设置
\texttt{mathescape=true}，或使用 \texttt{escapeinside} 指定转义标记。

\begin{codeblock}[language=SQL,caption={按客户汇总交易金额},label={lst:sql-summary},captionpos=b]
SELECT customer_id,
       COUNT(*) AS transaction_count,
       SUM(amount) AS total_amount
FROM transactions
GROUP BY customer_id
HAVING COUNT(*) >= 3
ORDER BY total_amount DESC;
\end{codeblock}

代码也可以保存在独立文件中，再用以下命令引入；其中路径应替换为自己的源文件：

\begin{codeblock}[language={[LaTeX]TeX}]
\lstinputlisting[style=lstStyleCode,language=Python,
  caption={程序标题},label={lst:my-program}]{path/to/program.py}
\end{codeblock}

\section{TikZ 流程图}

TikZ 可直接绘制矢量图。图\ref{fig:tikz-training-flow} 使用相对定位、判定节点和回路，
图中文字与论文一起排版。节点间距由 \texttt{positioning} 库控制，箭头由
\texttt{arrows.meta} 库提供。

\begin{figure}[htbp]
  \centering
  \begin{tikzpicture}[
    >=Latex,
    node distance=8mm and 13mm,
    processnode/.style={draw,rounded corners=2pt,fill=blue!5,
      minimum width=30mm,minimum height=9mm,align=center},
    decisionnode/.style={diamond,draw,aspect=2.5,inner sep=2pt,align=center},
    flow/.style={->,semithick}
  ]
    \node[processnode] (input) {读取并预处理数据};
    \node[processnode,below=of input] (fit) {计算目标函数};
    \node[decisionnode,below=of fit] (check) {满足停止条件？};
    \node[processnode,right=of check] (output) {输出模型};
    \node[processnode,below=of check] (update) {更新模型参数};
    \draw[flow] (input) -- (fit);
    \draw[flow] (fit) -- (check);
    \draw[flow] (check) -- node[above] {是} (output);
    \draw[flow] (check) -- node[right] {否} (update);
    \draw[flow] (update.west) -- ++(-12mm,0) |- (fit.west);
  \end{tikzpicture}
  \caption{迭代训练流程的 TikZ 示例}
  \label{fig:tikz-training-flow}
\end{figure}

\section{PGFPlots 数据曲线}

PGFPlots 在 TikZ 基础上提供坐标轴、图例和数据曲线。
图\ref{fig:pgfplots-loss} 直接使用示例坐标点，分别通过线型和标记区分两种方法，
便于在黑白打印时辨认。
实际使用时也可以用 \texttt{\textbackslash addplot table} 从数据文件读入数值。

\begin{figure}[htbp]
  \centering
  \begin{tikzpicture}
    \begin{axis}[
      width=0.88\linewidth,height=5.2cm,
      xlabel={迭代次数},ylabel={目标函数值},
      xmin=1,xmax=6,ymin=0,ymax=1.05,
      xtick={1,2,3,4,5,6},
      grid=major,grid style={gray!20},
      legend pos=north east,legend cell align=left,
      legend style={draw=none,fill=white,font=\zihao{5}},
      tick label style={font=\zihao{5}},
      label style={font=\zihao{5}}
    ]
      \addplot[blue!65!black,thick,mark=*]
        coordinates {(1,0.95) (2,0.65) (3,0.44) (4,0.31) (5,0.24) (6,0.20)};
      \addlegendentry{方法 A}
      \addplot[black!65,thick,dashed,mark=square*]
        coordinates {(1,0.95) (2,0.76) (3,0.61) (4,0.50) (5,0.43) (6,0.38)};
      \addlegendentry{方法 B}
    \end{axis}
  \end{tikzpicture}
  \caption{示例数据的目标函数变化曲线}
  \label{fig:pgfplots-loss}
\end{figure}

\section{数学环境、数值与单位}

\subsection{定义、定理与证明}

下面用样本均值展示定义、定理与证明的交叉引用。

\begin{definition}[样本均值]
  \label{def:sample-mean}
  对于 $n\ge1$ 个实数 $x_1,\ldots,x_n$，样本均值定义为
  \begin{equation}
    \bar{x}=\frac{1}{n}\sum_{i=1}^{n}x_i.
    \label{eq:sample-mean}
  \end{equation}
\end{definition}

\begin{theorem}[平方偏差最小化]
  \label{thm:least-squares-mean}
  设 $\bar{x}$ 如定义\ref{def:sample-mean} 所示。
  对任意实数 $a$，均有
  \begin{equation}
    \sum_{i=1}^{n}(x_i-a)^2
      \ge \sum_{i=1}^{n}(x_i-\bar{x})^2,
    \label{eq:least-squares-mean}
  \end{equation}
  等号当且仅当 $a=\bar{x}$ 时成立。
\end{theorem}

\begin{proof}
  根据式\eqref{eq:sample-mean}，有 $\sum_{i=1}^{n}(x_i-\bar{x})=0$，因此
  \begin{align*}
    \sum_{i=1}^{n}(x_i-a)^2
      &= \sum_{i=1}^{n}\bigl[(x_i-\bar{x})+(\bar{x}-a)\bigr]^2 \\
      &= \sum_{i=1}^{n}(x_i-\bar{x})^2+n(\bar{x}-a)^2.
  \end{align*}
  第二项非负，且仅当 $a=\bar{x}$ 时为零，故定理\ref{thm:least-squares-mean} 成立。
\end{proof}

\subsection{矩阵与分段函数}

% bmatrix 表示矩阵，cases 表示分段定义。
\begin{equation}
  \boldsymbol{X}=\begin{bmatrix}
    1 & x_1 \\
    \vdots & \vdots \\
    1 & x_n
  \end{bmatrix},
  \qquad
  h(z)=\begin{cases}
    z, & z\ge0, \\
    0, & z<0.
  \end{cases}
  \label{eq:matrix-and-cases}
\end{equation}

\subsection{数值对齐与单位}

\texttt{siunitx} 统一数值与单位的写法，例如
\num{1234567.89}、\qty{2.5}{\milli\second} 和 \qty{95.2}{\percent}。
表\ref{tab:numeric-results} 使用 \texttt{S} 列按小数点对齐，
表注由 \texttt{threeparttable} 提供。

\begin{table}[htbp]
  \centering
  \begin{threeparttable}
    \caption{小数点对齐的结果表格示例}
    \label{tab:numeric-results}
    \begin{tabular}{l S[table-format=2.2] S[table-format=3.1]}
      \toprule
      方法 & {准确率（\%）} & {耗时（\unit{\milli\second}）} \\
      \midrule
      基线方法 & 89.25 & 102.4 \\
      方法 A & 92.80 & 86.7 \\
      方法 B\tnote{a} & 94.15 & 91.2 \\
      \bottomrule
    \end{tabular}
    \begin{tablenotes}
      \item[a] 各列仅展示数值对齐和表注排版方式。
    \end{tablenotes}
  \end{threeparttable}
\end{table}
